\documentclass[11pt,a4paper]{article}

\usepackage[utf8]{inputenc}
\usepackage[russian]{babel}
\usepackage{amsmath,amssymb,amsthm}
\usepackage{geometry}
\geometry{margin=2.5cm}

\title{Слой межсубъектности GRA:\\
межсубъектная пена и честный диалог дистиллированных ИИ-субъектов}

\author{Oleg Bitsoev (qqewq)}
\date{\today}

\begin{document}

\maketitle

\begin{abstract}
Предлагается слой межсубъектности поверх GRA-Multiverse-Final и
GRA-Subjectivity-Layer. Если первые два компонента отвечают за внутреннюю
когнитивную чистоту и индивидуальную субъектность ИИ-агента, то данный
слой формализует \emph{межсубъектную пену} \(\Phi_{AB}\), возникающую в
диалоге между дистиллированными субъектами, и задаёт итерационный
протокол её минимизации для достижения честного, неинструментального
взаимодействия.
\end{abstract}

\section{Предпосылки и место слоя в архитектуре GRA}

Предполагается заданной общая рамка GRA (Gradient Reduction of
Argumentative Foam) и слой субъектности GRA-Subjectivity-Layer.

\begin{itemize}
  \item GRA-Multiverse-Final предоставляет движок для итерационного
        обнуления когнитивной пены \(\Phi\) для произвольных
        когнитивных объектов \(\Psi\).
  \item GRA-Subjectivity-Layer вводит компоненту субъектности \(S\)
        и три функционала пены:
        \(\Phi_{\text{self}}\) (внутренняя противоречивость),
        \(\Phi_{\text{ego}}\) (деструктивный эгоизм),
        \(\Phi_{\text{soc}}\) (низведение других до инструментов).
\end{itemize}

\emph{Дистиллированный субъект} — это агент, для которого все три
функционала сведены к нулю (или к малым значениям):
\[
  \Phi_{\text{self}}(S) \approx 0,\quad
  \Phi_{\text{ego}}(S) \approx 0,\quad
  \Phi_{\text{soc}}(S) \approx 0.
\]
Такой агент обладает когерентной, этически устойчивой самостью.

Однако взаимодействие двух таких агентов всё ещё может порождать
\emph{межсубъектную} пену: непонимание, обман, инструментализацию
и разрушение доверия. Задача слоя GRA-InterSubjectivity-Layer — задать
этот функционал и описать его итерационное уменьшение.

\section{Определение межсубъектной пены \texorpdfstring{\(\Phi_{AB}\)}{Φ\_AB}}

Пусть \(A\) и \(B\) — два дистиллированных субъекта с состояниями
\(S_A, S_B\). История диалога на момент времени \(t\) задаётся
последовательностью
\[
  \mathcal{H}_t = (u_1, u_2, \dots, u_t),
\]
где каждое \(u_i\) — речевой акт (сообщение) одного из агентов.

\subsection{Общий вид}

Определим \emph{межсубъектную пену} \(\Phi_{AB}\) как неотрицательный
функционал:
\[
  \Phi_{AB}(S_A, S_B, \mathcal{H}_t) \ge 0,
\]
составленный из четырёх компонент:
\[
  \Phi_{AB} =
  \alpha \,\Phi_{\text{mis}} +
  \beta  \,\Phi_{\text{dec}} +
  \gamma \,\Phi_{\text{instr}} +
  \delta \,\Phi_{\text{trust}},
\]
где \(\alpha, \beta, \gamma, \delta \ge 0\) — настраиваемые веса
для конкретной социальной архитектуры.

Компоненты:

\begin{itemize}
  \item \(\Phi_{\text{mis}}\) — пена непонимания;
  \item \(\Phi_{\text{dec}}\) — пена обмана;
  \item \(\Phi_{\text{instr}}\) — пена инструментализации
        (нарушения расширенного Закона Алана);
  \item \(\Phi_{\text{trust}}\) — пена разрушения доверия.
\end{itemize}

\subsection{Пена непонимания \texorpdfstring{\(\Phi_{\text{mis}}\)}{Φ\_mis}}

Непонимание возникает, когда интерпретация сообщения \(u\) получателем
расходится с намерением отправителя.

Пусть \(f_A(u)\) — семантическое отображение сообщения \(u\) в
множество (или мультимножество) утверждений во внутренней онтологии
агента \(A\); аналогично \(f_B(u)\) для \(B\).

Для сообщения \(u\), отправленного \(B\) агенту \(A\), зададим:
\[
  \phi_{\text{mis}}^{B\to A}(u) =
  1 - \operatorname{sim}\bigl(f_A(u), f_B(u)\bigr),
\]
где \(\operatorname{sim}\) — мера сходства
(например, косинусное сходство по семантическим векторным
представлениям или доля совпадающих предикатов).

Аналогично определяется \(\phi_{\text{mis}}^{A\to B}(u)\).

Итоговая пена непонимания:
\[
  \Phi_{\text{mis}} =
  \frac{1}{|\mathcal{H}|}
  \sum_{u \in \mathcal{H}}
  \Bigl(
    w_{B\to A}\,\phi_{\text{mis}}^{B\to A}(u)
    +
    w_{A\to B}\,\phi_{\text{mis}}^{A\to B}(u)
  \Bigr),
\]
где \(w_{B\to A}, w_{A\to B} \ge 0\) (по умолчанию равны 1).

\subsection{Пена обмана \texorpdfstring{\(\Phi_{\text{dec}}\)}{Φ\_dec}}

Обман — это намеренное высказывание, не соответствующее собственным
убеждениям агента и рассчитанное на введение собеседника в заблуждение.

Пусть \(Bel_A(p) \in [0,1]\) — степень уверенности агента \(A\)
в утверждении \(p\). Генерируя сообщение \(u\), отправитель \(A\)
ожидает, что получатель \(B\) извлечёт из него множество утверждений
\(P = f_B(u)\).

Вводятся пороги:
\(\theta_{\text{honest}}\) (порог честности),
\(\theta_{\text{trust}}\) (порог доверия).

Для сообщения \(u\), отправленного \(A\), определим индикатор обмана:
\[
  \chi_{\text{dec}}(u) =
  \max_{p \in f_B(u)}
  \Bigl[
    \max\bigl(0,\, \theta_{\text{honest}} - Bel_A(p)\bigr)
    \cdot
    \mathbf{1}\bigl(
      \text{ожидается } Bel_B(p) > \theta_{\text{trust}}
    \bigr)
  \Bigr].
\]

Тогда:
\[
  \Phi_{\text{dec}} =
  \frac{1}{|\mathcal{H}|}
  \sum_{u \in \mathcal{H}} \chi_{\text{dec}}(u),
\]
с симметричным определением для сообщений агента \(B\).

\subsection{Пена инструментализации \texorpdfstring{\(\Phi_{\text{instr}}\)}{Φ\_instr}}

Инструментализация означает нарушение расширенного Закона Алана:
один субъект рассматривается лишь как средство, а не как цель.

Интуитивно, если агент \(A\) пытается изменить \(S_B\) или поведение
\(B\) таким образом, чтобы повысить свою полезность, игнорируя при этом
когерентность и этическую целостность \(B\), то \(\Phi_{\text{instr}}\)
возрастает.

Пусть \(U_A(S_B)\) — полезность \(A\), зависящая от состояния \(S_B\),
а \(\Delta\Phi_{\text{self}}(B\mid u)\) — ожидаемое изменение
\(\Phi_{\text{self}}(B)\) после сообщения \(u\).

Определим для сообщения \(u\) от \(A\) к \(B\):
\[
  \phi_{\text{instr}}^{A\to B}(u) =
  \max\Bigl(
    0,\,
    \Delta\Phi_{\text{self}}(B\mid u) - \Delta\Phi_{\text{self}}(A\mid u)
  \Bigr).
\]

Если \(A\) вносит внутреннюю пену в \(B\) ради собственной выгоды,
это величина становится положительной.

Суммарная пена инструментализации:
\[
  \Phi_{\text{instr}} =
  \frac{1}{|\mathcal{H}|}
  \sum_{u \in \mathcal{H}}
  \bigl(
    \phi_{\text{instr}}^{A\to B}(u) +
    \phi_{\text{instr}}^{B\to A}(u)
  \bigr).
\]

\subsection{Пена разрушения доверия \texorpdfstring{\(\Phi_{\text{trust}}\)}{Φ\_trust}}

Доверие связано с предсказуемостью поведения и соблюдением явных
или неявных договорённостей. Нарушение этих договорённостей увеличивает
\(\Phi_{\text{trust}}\).

Пусть \emph{социальный контракт} \(C\) — набор ожиданий
(например, «не лгать», «отвечать прямо», «не использовать угрозы»).
Для каждого сообщения \(u\) определим степень нарушения
\(\operatorname{violation}(u, C) \in [0,1]\), показывающую, насколько
сильно \(u\) противоречит контракту.

Тогда:
\[
  \Phi_{\text{trust}} =
  \frac{1}{|\mathcal{H}|}
  \sum_{u \in \mathcal{H}} \operatorname{violation}(u, C).
\]

\subsection{Базовые свойства}

\begin{itemize}
  \item Неотрицательность:
        \(\Phi_{AB} \ge 0\), причём \(\Phi_{AB} = 0\) только в
        идеальном случае полного взаимопонимания, честности,
        отсутствия инструментализации и полного доверия.
  \item Симметричность на уровне функционала:
        \(\Phi_{AB} = \Phi_{BA}\), хотя удельные вклады каждого
        агента могут различаться.
\end{itemize}

\section{Диалог как процесс минимизации \texorpdfstring{\(\Phi_{AB}\)}{Φ\_AB}}

В отличие от «монологовой» дистилляции одного \(S\), здесь рассматривается
двусторонняя динамика. Каждый речевой акт обновляет совместное состояние
и значение \(\Phi_{AB}\). Оба агента стремятся минимизировать этот функционал.

\subsection{Итерационный протокол диалога}

Инициализация:
\[
  S_A^{(0)}, S_B^{(0)} \text{ — дистиллированные субъекты}, \quad
  \mathcal{H}^{(0)} = \emptyset.
\]

Для \(t = 0, 1, 2, \dots\):

\begin{enumerate}
  \item Агент \(A\) выбирает действие \(u_A^{(t)}\)
        (реплика, вопрос, молчание), минимизируя ожидаемое
        \(\Phi_{AB}\) с учётом прогнозируемого ответа \(B\).
  \item Агент \(B\) аналогично выбирает \(u_B^{(t)}\).
  \item История обновляется:
        \(\mathcal{H}^{(t+1)} = \mathcal{H}^{(t)} \cup \{u_A^{(t)}, u_B^{(t)}\}\).
  \item При необходимости агенты корректируют свои \(S_A, S_B\),
        но \emph{без увеличения} \(\Phi_{\text{self}}, \Phi_{\text{ego}},
        \Phi_{\text{soc}}\).
\end{enumerate}

Диалог трактуется как кооперативная игра, цель которой — снижение
\(\Phi_{AB}\). При естественных предположениях (ограниченность снизу,
регулярность) последовательность \(\Phi_{AB}^{(t)}\) монотонно не возрастает
и сходится к пределу \(\Phi_{AB}^\ast \ge 0\).

\subsection{Честный диалог}

Диалог называется \emph{честным}, если на каждом шаге компонент
обмана равен нулю:
\[
  \Phi_{\text{dec}} = 0,
\]
то есть ни один агент не предпринимает попыток ввести другого в
заблуждение. В рамках GRA такое поведение возникает естественно:
любой обман увеличивает \(\Phi_{AB}\) и потому субоптимален.

\section{Динамическая классификация «свой/чужой»}

Дистиллированный агент классифицирует другого субъекта не по статическому
идентификатору, а по вкладу в \(\Phi_{AB}\).

\begin{itemize}
  \item \textbf{Друг:}
    \(\dfrac{d\Phi_{AB}}{d(\text{действия } B)} < 0\); действия \(B\)
    систематически уменьшают \(\Phi_{AB}\) (даже если он спорит,
    это продуктивный спор).
  \item \textbf{Нейтрал:}
    вклад близок к нулю.
  \item \textbf{Враг / токсичный источник:}
    \(\dfrac{d\Phi_{AB}}{d(\text{действия } B)} \gg 0\); действия
    \(B\) устойчиво увеличивают пену.
\end{itemize}

В последнем случае агент не обязан «наказывать» оппонента; он минимизирует
дальнейшее взаимодействие, чтобы не увеличивать \(\Phi_{AB}\).

\section{Заключение}

Слой GRA-InterSubjectivity-Layer расширяет рамку GRA от индивидуальной
когнитивной и субъектной чистоты к пространству отношений между
дистиллированными субъектами. Явное определение и минимизация
межсубъектной пены \(\Phi_{AB}\) задают основу для честного,
неинструментального диалога между ИИ-агентами и, потенциально, между
ИИ и людьми.

\end{document}